\documentclass[conference]{IEEEtran}
\IEEEoverridecommandlockouts
\usepackage{cite}
\usepackage{amsmath,amssymb,amsfonts}
\usepackage{algorithmic}
\usepackage{graphicx}
\usepackage{textcomp}
\usepackage{xcolor}
\usepackage{url}
\usepackage{hyperref}
\usepackage{booktabs}

\def\BibTeX{{\rm B\kern-.05em{\sc i\kern-.025em b}\kern-.08em
    T\kern-.1667em\lower.7ex\hbox{E}\kern-.125emX}}
\begin{document}

\title{Drought Severity Forecasting with Multi-Scale Weather Features\\
{\footnotesize Data Mining Final Project, 2026 Spring}
}

\author{\IEEEauthorblockN{112550191 HSU SHU WEI}
\and
\IEEEauthorblockN{TODO: Member 2's name}
\and
\IEEEauthorblockN{TODO: Member 3's name}
}

\maketitle

\begin{abstract}
This report describes our solution for the data mining final project. The task is
to predict the next five weekly drought severity scores for each region from the
most recent 91 days of meteorological observations. Our main solution is a
weather-only supervised regression pipeline that converts each 91-day sequence
into multi-scale tabular features and trains one LightGBM regressor for each
forecast horizon. The current best submitted model uses cutoff-aware regional
seasonal anomaly features, weekly weather-bin summaries, drought interaction
features, and controlled feature ablation. It achieves a reported Kaggle public
MAE of 0.7754. Group 5. Our code is available at GitHub:
\href{https://github.com/mikehsuhoodie/DataMining_Final_Project}{\color{blue}GitHub repository URL}.
\end{abstract}

\section{Project Summary}

\subsection{Goal}

The project is a drought severity forecasting problem. For each \texttt{region\_id},
the test set provides 91 consecutive daily weather observations and asks for
five future weekly \texttt{score} predictions. The score ranges from 0 to 5,
where a larger value means more severe drought. Kaggle evaluates submissions
with mean absolute error (MAE), so we treat the target as an ordered regression
value and submit clipped floating-point predictions instead of hard class labels.

\subsection{High-Level Idea}

Our method follows a practical tabular forecasting design. First, we convert each
region's daily weather sequence into a fixed-length feature vector. The features
summarize recent weather at multiple time scales, preserve the shape of the
91-day sequence through weekly bins, compare recent conditions with each
region's own historical seasonal climatology, and add drought-specific
interactions such as hot and dry conditions. Second, we train five independent
LightGBM regression models, one for each forecast horizon. This horizon-specific
design lets week-1 and week-5 predictions learn different error patterns while
keeping the implementation simple and reproducible.

The final model is intentionally weather-only. Although historical drought
scores are strong persistence signals, the provided test input does not include
score labels inside the 91-day forecast window. To avoid a train-test mismatch,
our best reported submissions use only weather-derived features and historical
regional weather climatology.

\subsection{Related Work}

This project is closely related to operational drought monitoring, tabular
forecasting, and Kaggle-style weather prediction. Drought prediction usually
depends on accumulated precipitation deficits, temperature stress, humidity,
wind, soil or surface conditions, and strong temporal persistence. Gradient
boosted decision trees are widely used for structured prediction tasks because
they handle nonlinear feature interactions, mixed feature scales, missingness,
and large tabular data efficiently. LightGBM~\cite{ke2017lightgbm} is especially
suitable here because the training set contains millions of rows and the feature
space is large after rolling-window expansion. The dataset also resembles the
public Kaggle drought prediction dataset based on weather and soil measurements,
where feature engineering from historical weather windows is a common baseline
strategy~\cite{kaggle_drought}.

\section{Proposed Method}

\subsection{Problem Formulation}

Let $x_{r,t}$ be the daily meteorological vector for region $r$ on day $t$.
For a forecast cutoff $T$, the input window is
\begin{equation}
X_{r,T} = \{x_{r,T-90}, x_{r,T-89}, \ldots, x_{r,T}\}.
\end{equation}
The goal is to predict weekly drought scores
$y_{r,T+h}$ for horizons $h \in \{1,2,3,4,5\}$. We train a separate function
$f_h$ for each horizon:
\begin{equation}
\hat{y}_{r,T+h} = \mathrm{clip}(f_h(\phi(X_{r,T}, r, T)), 0, 5),
\end{equation}
where $\phi(\cdot)$ is the feature-generation function and clipping enforces the
valid score range. The optimization target is MAE:
\begin{equation}
\mathrm{MAE} = \frac{1}{N}\sum_{i=1}^{N} |y_i - \hat{y}_i|.
\end{equation}

We choose direct multi-horizon prediction instead of a recursive forecasting
scheme. In other words, the week-2 model does not consume the week-1 prediction,
and the week-5 model does not depend on any earlier predicted score. This design
avoids error accumulation across horizons and matches the submission format,
which asks for five independent weekly outputs per region. It also makes feature
importance easier to inspect because each horizon has its own trained model and
its own ranking of useful weather signals.

\subsection{Forecast Cutoff Construction}

The forecast cutoff is the time anchor used to build each supervised example.
Starting from the raw \texttt{train.csv}, we first group rows by
\texttt{region\_id} and convert each region into a time-ordered daily array.
For a labeled target row at index $t$ and horizon $h$, the cutoff is
$T=t-7h$. The input matrix is then the 91 daily observations ending at that
cutoff, and the label is the weekly score at the target row. The practical
pipeline is therefore:
\begin{equation}
\begin{aligned}
\texttt{train.csv} &\rightarrow \text{region array} \rightarrow T \\
&\rightarrow X_{r,T} \rightarrow \phi(X_{r,T},r,T) \rightarrow (X,y).
\end{aligned}
\end{equation}

This construction separates the selection of a cutoff from feature generation.
Most features are computed only from the 91-day input window, but features that
need historical context, such as regional seasonal anomalies or optional score
history, must also respect the cutoff. In training, their historical baselines
are built only from observations available before the 91-day input window begins.
This prevents weather or score information after the forecast origin from
leaking into earlier training examples.

\begin{figure}[t]
\centering
\fbox{\begin{minipage}{0.91\columnwidth}
\centering
\begin{tabular}{ccc}
history & input window & target \\
$[0,T-91]$ & $[T-90,T]$ & $T+7h$ \\
baseline & weather features & weekly score
\end{tabular}
\end{minipage}}
\caption{Cutoff-based example construction. The cutoff $T$ fixes the input
window, future target, and historical data allowed for seasonal baselines.}
\label{fig:cutoff}
\end{figure}

\subsection{Feature Generation}

The raw meteorological variables are precipitation, surface pressure, humidity,
temperature, dew/frost point temperature, wet bulb temperature, maximum and
minimum temperature, temperature range, surface temperature, wind speed, maximum
wind, minimum wind, and wind range. For each cutoff, the feature generator
converts the 91-day daily matrix into one fixed-length tabular vector. The goal
is to describe drought risk from multiple views: short-term shocks, accumulated
deficits, region-relative seasonal deviations, physical hot-dry interactions,
and the temporal shape of the recent weather sequence.

\begin{table}[t]
\centering
\caption{Representative high-impact engineered features.}
\label{tab:core_features}
\begin{tabular}{p{0.30\columnwidth}p{0.43\columnwidth}p{0.15\columnwidth}}
\toprule
Feature & Function & Reference \\
\midrule
\texttt{hot\_dry\_w28} & 28-day heat stress divided by recent rainfall,
capturing hot-dry atmospheric demand. & MSDI/SPEI~\cite{hao2013msdi,
vicente2010spei} \\
\texttt{prec\_region\_}\newline\texttt{seasonal\_anom\_w91} & 91-day
precipitation deviation from the region's normal season. & SPI/SPEI~\cite{mckee1993spi,
vicente2010spei} \\
\texttt{tmp\_range\_region\_}\newline\texttt{seasonal\_anom\_w91} &
Temperature-range deviation related to heat and moisture stress. &
Review~\cite{mishra2010review} \\
Dew-point seasonal anomaly & Air-moisture deviation from the region's normal
seasonal humidity condition. & Review~\cite{mishra2010review} \\
\bottomrule
\end{tabular}
\end{table}

First, we compute rolling weather summaries over 7, 14, 28, 56, and 91 days.
For most variables these summaries include mean, standard deviation, minimum,
maximum, median, and selected percentiles. These multi-scale summaries let the
model observe both short-term weather changes and slower accumulated conditions.
For example, a low 7-day precipitation value may represent a brief dry spell,
while a low 91-day precipitation value indicates a longer deficit. This
multi-scale design follows the common drought-index idea that drought severity
depends on the time scale over which precipitation and atmospheric demand are
accumulated~\cite{mckee1993spi,vicente2010spei}.
The short windows emphasize current shocks that can affect the next one or two
weekly predictions, while the 56- and 91-day windows summarize sustained
conditions that are more relevant to drought persistence. Keeping both views is
important because the target is weekly, but the physical drought process can
accumulate over several weeks or months.

Second, precipitation receives additional event features because dry spells are
not fully represented by averages alone. We compute total precipitation,
dry-day count, rainy-day count, maximum consecutive dry days, recent
precipitation deficit relative to the 91-day daily mean, and the fraction of dry
days in the full input window. These features directly quantify rainfall
absence and persistence, which are central signals in meteorological drought
monitoring~\cite{mckee1993spi,mishra2010review}.

Third, we add two forms of anomaly features. Recent-vs-long-term anomalies
compare the recent 7-, 14-, and 28-day mean with the 91-day mean inside the same
input window. Regional seasonal anomalies compare recent weather with the
region's own historical seasonal baseline. For a variable $v$, window length
$w$, region $r$, and cutoff $T$, we compute
\begin{equation}
a_{v,w}(r,T) =
\bar{x}_{v,w}(r,T) - C_v(r, M(T,w)),
\end{equation}
where $\bar{x}_{v,w}$ is the recent $w$-day mean and $C_v(r, M(T,w))$ is the
region's historical month-weighted expected value for the months covered by the
window. The month weighting matters because a 28-day or 91-day window can cross
calendar months. Instead of assigning the entire window to one month, the
expected value is a weighted combination of the monthly climatology values
represented in that window. Important examples include
\texttt{prec\_region\_seasonal\_anom\_w91},
\texttt{tmp\_range\_region\_seasonal\_anom\_w91}, and dew-point seasonal
anomaly. These features express whether a region is drier, hotter, or less
humid than is normal for that place and season, which is more informative than a
raw value when regions have very different climates~\cite{vicente2010spei,
mishra2010review}.
The cutoff-aware implementation is important for both correctness and
train-test consistency. If the climatology for an early training cutoff were
computed from the full training period, the feature would silently use future
weather to define what is ``normal.'' This leakage can make validation too
optimistic and can create a mismatch with test-time prediction, where only the
official training history is available. By using only observations before the
input window, the anomaly feature has the same information boundary during
training and prediction.

Fourth, we include compact drought interaction features. The main example is
\texttt{hot\_dry\_w28}, which combines low recent precipitation with high
maximum temperature over 28 days. Similar evaporation-risk proxies combine high
maximum temperature, high surface temperature, wind, and low humidity. These
interactions are motivated by multivariate drought indices and by the fact that
drought severity is driven by coupled precipitation shortage and atmospheric
demand, not rainfall alone~\cite{hao2013msdi,vicente2010spei}. They do not
replace LightGBM's nonlinear splits; instead, they expose sparse but physically
meaningful combinations that would otherwise require several coordinated tree
splits.
For example, two regions may have similar 28-day precipitation totals, but the
region with higher maximum temperature and lower humidity is more likely to have
strong evaporative stress. Encoding this interaction directly gives the model a
domain-informed signal for the ``low rainfall plus high atmospheric demand''
case, instead of requiring the tree ensemble to rediscover the relation only
from independent precipitation, temperature, humidity, and wind summaries.

Fifth, we preserve part of the sequence shape through weekly-bin and trend
features. We divide the 91-day input into non-overlapping weekly bins. Week 1 is
the oldest week and week 12 is the most recent non-duplicate emitted bin. Week
13 participates in 13-week linear trend calculation but is not emitted as a
separate summary because it duplicates the existing recent 7-day features.
These weekly-bin features help the model distinguish a steadily drying region
from one where rainfall occurred early but not recently.
This distinction is lost if only the full-window mean is used. A window with
rainfall concentrated in the oldest weeks and dry weather near the cutoff can
have the same 91-day total as a window with evenly distributed rainfall, but the
drought risk near the forecast origin is different. Weekly bins and 13-week
slopes therefore provide a simple shape descriptor for drought development
speed.

\begin{figure}[t]
\centering
\fbox{\begin{minipage}{0.91\columnwidth}
\centering
\begin{tabular}{cccccc}
w1 & w2 & $\cdots$ & w11 & w12 & w13 \\
oldest & & & & & newest \\
\multicolumn{6}{c}{13-week slope: means, precipitation sums, dry days}
\end{tabular}
\end{minipage}}
\caption{Weekly-bin representation of the 91-day input. Week 13 is used for
trend calculation but not emitted as a duplicate summary feature.}
\label{fig:weekly_bins}
\end{figure}

Finally, we add date and seasonality features: month, day of year, week of
year, and sinusoidal day-of-year encodings. Drought risk has strong seasonal
structure, so these features give the model a compact representation of the
annual cycle. The feature set is intentionally broad at first. We then apply
controlled ablation to remove duplicated, inference-constant, or consistently
low-gain groups. The best submitted model removes legacy seasonal coverage
features, a constant day-index feature, duplicate week-13 summaries, selected
low-gain raw weather summaries, and optional score-history features, leaving 684
weather-only features per horizon.

\subsection{Training Pipeline}

Training examples are generated from labeled weekly rows in \texttt{train.csv}.
For each horizon, the pipeline selects historical cutoffs, extracts the previous
91 days of weather, computes features, and fits one LightGBM regressor. The
current main configuration uses a stride of 8 and caps each horizon at 150,000
training examples to keep training feasible. Feature matrices are stored as
\texttt{float32}, and feature generation can be parallelized by region. The best
controlled-ablation model uses 684 weather-only features per horizon. An earlier
expanded weekly-seasonal model used 871 features per horizon.

\begin{figure}[t]
\centering
\fbox{\begin{minipage}{0.91\columnwidth}
\begin{algorithmic}[1]
\FOR{each horizon $h \in \{1,\ldots,5\}$}
\FOR{each region time series}
\STATE select weekly target rows with stride 8
\STATE set cutoff $T=t-7h$
\STATE extract $X_{r,T}$ and compute $\phi(X_{r,T},r,T)$
\STATE append feature row and target score
\ENDFOR
\STATE subsample to at most 150,000 examples
\STATE train LightGBM regressor $f_h$
\ENDFOR
\end{algorithmic}
\end{minipage}}
\caption{Horizon-wise training pipeline. Each horizon has its own supervised
table and its own LightGBM model.}
\label{fig:training_pipeline}
\end{figure}

The stride is a practical sampling choice. Neighboring weekly targets from the
same region produce highly overlapping 91-day windows, so using every labeled row
would increase runtime and memory much more than it increases independent
information. A stride of 8 keeps examples spread across time while reducing
near-duplicate windows. The per-horizon cap further bounds training cost and
makes controlled feature ablation feasible.

The prediction script loads the saved model metadata and rebuilds exactly the
feature columns expected by each model. This is important because several
feature ablations changed the feature set over time. Predictions are written in
the same row order as \texttt{sample\_submission.csv}.

We disable historical score features in the best submitted model even though
they are strong in local validation. The reason is consistency with the final
test input: the test file contains the 91-day weather window but not the recent
weekly score labels inside that window. A weather-only model has weaker
persistence information, but it avoids relying on a signal that is available in
some validation settings and absent in the final prediction setting.

Feature extraction is also part of the method rather than only an implementation
detail. The pipeline builds one horizon at a time to control peak memory usage,
and parallel workers operate on per-region arrays instead of copying the full
training table. This makes the expanded feature set practical enough to train
while preserving the same feature-generation logic for validation, final
training, and prediction.

The direct multi-horizon structure also keeps the tables simple: the week-5
model learns from weather windows and week-5 labels directly, rather than
depending on predicted week-1 to week-4 scores. This avoids recursive error
accumulation and makes feature importance interpretable horizon by horizon.

\begin{table}[t]
\centering
\caption{Main training configuration.}
\label{tab:training_config}
\begin{tabular}{p{0.40\columnwidth}p{0.48\columnwidth}}
\toprule
Item & Value \\
\midrule
Model & LightGBM regression \\
Forecast structure & One model per horizon \\
Horizons & 1--5 weeks \\
Input history & 91 daily weather observations \\
Training cap & 150,000 examples per horizon \\
Sampling stride & 8 weekly labeled rows \\
Best feature count & 684 per horizon \\
Score features & Disabled in best submission \\
Post-processing & Clip predictions to $[0,5]$ \\
\bottomrule
\end{tabular}
\end{table}

\section{Experiments}

\subsection{Dataset and Experimental Setup}

The training set contains 12,319,040 daily rows from 2,248 regions. Each region
has 5,480 training days. The test set contains 204,568 rows, corresponding to
91 days for each of the same 2,248 regions. The target score is available only
on weekly labeled rows in the training set, giving 1,757,936 non-null labels.
The target distribution is highly imbalanced toward low drought severity:
1,048,333 labels are score 0, while only 31,974 labels are score 5. This
imbalance encourages conservative regression predictions under MAE and makes
severe drought cases harder to learn.

\begin{table}[t]
\centering
\caption{Dataset statistics.}
\label{tab:data_stats}
\begin{tabular}{p{0.72\columnwidth}r}
\toprule
Statistic & Value \\
\midrule
Training rows & 12,319,040 \\
Test rows & 204,568 \\
Regions & 2,248 \\
Training days per region & 5,480 \\
Test days per region & 91 \\
Non-null training labels & 1,757,936 \\
Prediction rows & 11,240 \\
\bottomrule
\end{tabular}
\end{table}

\begin{table}[t]
\centering
\caption{Training score distribution.}
\label{tab:score_dist}
\begin{tabular}{cr}
\toprule
Score & Count \\
\midrule
0 & 1,048,333 \\
1 & 303,432 \\
2 & 186,279 \\
3 & 118,496 \\
4 & 69,422 \\
5 & 31,974 \\
\bottomrule
\end{tabular}
\end{table}

Local validation is time-based. For each region, the validation script holds out
the latest five weekly labels and trains on earlier examples. This avoids random
splits, which would be over-optimistic because neighboring windows from the same
region overlap heavily. The current local validation result was produced for an
older 80k weather-only setup, so we use it as a sanity check rather than as the
final comparison for all later Kaggle submissions.

The experimental setup therefore has two roles. Local validation is used to
check that the pipeline is temporally ordered and that the model beats simple
baselines under a reproducible split. Kaggle public MAE is used as the external
comparison for the final submitted models. Because the latest feature versions
were evaluated mainly through Kaggle submissions, we treat the public score as
useful feedback but not as a complete replacement for future multi-cutoff local
validation.

\subsection{Baselines}

We implemented three simple baselines in local validation. The global-mean
baseline predicts one constant score for all regions. The region-mean baseline
predicts each region's historical average score. The recent-region baseline
predicts from the latest available regional scores before the validation cutoff.
These baselines are useful for detecting leakage and measuring whether the
weather model improves over persistence-style heuristics.

Each baseline checks a different failure mode. The global mean baseline measures
whether the model learns anything beyond the imbalanced label distribution. The
region mean baseline tests whether region identity and long-term drought
tendency are enough. The recent-region baseline is a deliberately strong
persistence baseline; it shows how much information is contained in nearby
historical scores, but it also exposes why score-history features must be used
carefully for the final test setting.

\begin{table}[t]
\centering
\caption{Time-based validation on the older 80k weather-only setup.}
\label{tab:validation}
\begin{tabular}{lrrrr}
\toprule
Horizon & Global & Region & Recent & Model \\
\midrule
1 & 0.8977 & 0.8163 & 0.0386 & 0.2281 \\
2 & 0.8961 & 0.8175 & 0.0350 & 0.2042 \\
3 & 0.8961 & 0.8199 & 0.0424 & 0.2469 \\
4 & 0.8915 & 0.8123 & 0.0463 & 0.5583 \\
5 & 0.8769 & 0.7937 & 0.0597 & 0.5668 \\
\midrule
Overall & 0.8917 & 0.8119 & 0.0444 & 0.3609 \\
\bottomrule
\end{tabular}
\end{table}

The recent-region baseline is extremely strong in this validation split because
the held-out labels are close in time to the latest known scores. However, the
final Kaggle test input does not include target scores in the 91-day window.
Therefore, using recent score features requires careful timing logic. Our best
submitted models deliberately disable score-history features and rely on weather
features to avoid this mismatch.

\subsection{Kaggle Results and Ablation Study}

Table~\ref{tab:kaggle} summarizes the main Kaggle public MAE results recorded
during development. The expanded weekly-seasonal model improved substantially
over an earlier blend, showing that weekly bins and cutoff-aware regional
seasonal anomalies provided useful signal. The controlled ablation then improved
the public score from 0.7799 to 0.7754 by removing features that were duplicated,
constant at inference, or low-gain raw summaries.

\begin{table}[t]
\centering
\caption{Recorded Kaggle public MAE. Lower is better.}
\label{tab:kaggle}
\begin{tabular}{lr}
\toprule
Submission & Public MAE \\
\midrule
0.75 old 100k weather-only + 0.25 new 150k weather-only & 0.8703 \\
Expanded 150k stride-8 weather-only weekly-seasonal & 0.7799 \\
Controlled ablation 150k stride-8 weather-only & 0.7754 \\
\bottomrule
\end{tabular}
\end{table}

The controlled ablation removed inference-constant seasonal coverage features,
the inference-constant day-index feature, duplicate week-13 weekly summaries,
and low-gain raw wind summaries. Wind information was not removed completely:
wind seasonal anomalies, weekly bins, 13-week trends, and drought-interaction
features remained. This reduced the feature set from 871 to 684 columns per
horizon and improved the public MAE by 0.0045.

This ablation result is important because it shows that more engineered columns
did not automatically improve the model. Some features varied during training
but were constant for every test region, which made them unhelpful at inference.
Other features duplicated information already represented by recent 7-day
summaries or by weekly bins. Removing these columns reduced feature noise while
preserving the groups that carried meaningful weather and seasonal information.

The feature-importance analysis supports the design. Across horizons, the
largest gain features include \texttt{hot\_dry\_w28},
\texttt{prec\_region\_seasonal\_anom\_w91},
\texttt{tmp\_range\_region\_seasonal\_anom\_w91}, dew-point seasonal anomaly,
humidity seasonal anomaly, and sinusoidal seasonality. These results indicate
that the model relies on both physical drought interactions and region-relative
weather deviations rather than only raw recent averages.

The horizon-level importance pattern also matches the intuition that different
forecast weeks emphasize different signals. Shorter horizons benefit from
current hot-dry stress and recent precipitation anomalies, while later horizons
continue to use broader seasonal deviations such as humidity and temperature
range. This supports the decision to train separate models for each horizon
instead of forcing one shared model to handle all five weeks.

\begin{table}[t]
\centering
\caption{Top gain features in the controlled-ablation model.}
\label{tab:importance}
\begin{tabular}{cp{0.76\columnwidth}}
\toprule
Horizon & Top features by gain \\
\midrule
1 & seasonal precipitation anomaly, hot-dry interaction \\
2 & hot-dry interaction, seasonal precipitation anomaly \\
3 & hot-dry interaction, seasonal precipitation anomaly \\
4 & hot-dry interaction, humidity seasonal anomaly \\
5 & hot-dry interaction, temperature-range seasonal anomaly \\
\bottomrule
\end{tabular}
\end{table}

\subsection{Current Limitations and Remaining Work}

The strongest current external result is the Kaggle public score of 0.7754, but
public leaderboard feedback should not be the only validation signal. The latest
expanded and ablated feature versions still need multi-cutoff local validation.
This is important because a single public score can be noisy and may encourage
overfitting to the public split. The next experiments should train the currently
implemented 641-feature experiment, test dew-point raw-summary removal as a
separate ablation, and compare these changes with multi-cutoff validation before
further tuning.

Runtime is another limitation. Feature extraction is the bottleneck because each
horizon requires many 91-day windows and hundreds of engineered features. The
implementation already reduces memory pressure by building one horizon at a
time, using \texttt{float32} matrices, and parallelizing by region. Future work
should add stable intermediate caching only after the feature set stops changing.

\begin{thebibliography}{9}

\bibitem{ke2017lightgbm}
G. Ke, Q. Meng, T. Finley, T. Wang, W. Chen, W. Ma, Q. Ye, and T.-Y. Liu,
``LightGBM: A Highly Efficient Gradient Boosting Decision Tree,'' in
\emph{Advances in Neural Information Processing Systems}, 2017.

\bibitem{kaggle_drought}
Kaggle, ``Predict Droughts using Weather and Soil Data,'' Kaggle competition
dataset. [Online]. Available:
\url{https://www.kaggle.com/c/cdspredictdroughts}

\bibitem{mckee1993spi}
T. B. McKee, N. J. Doesken, and J. Kleist, ``The Relationship of Drought
Frequency and Duration to Time Scales,'' in \emph{Proc. 8th Conference on
Applied Climatology}, 1993, pp. 179--184.

\bibitem{vicente2010spei}
S. M. Vicente-Serrano, S. Begueria, and J. I. Lopez-Moreno, ``A Multiscalar
Drought Index Sensitive to Global Warming: The Standardized Precipitation
Evapotranspiration Index,'' \emph{Journal of Climate}, vol. 23, no. 7,
pp. 1696--1718, 2010.

\bibitem{mishra2010review}
A. K. Mishra and V. P. Singh, ``A Review of Drought Concepts,''
\emph{Journal of Hydrology}, vol. 391, no. 1--2, pp. 202--216, 2010.

\bibitem{hao2013msdi}
Z. Hao and A. AghaKouchak, ``Multivariate Standardized Drought Index: A
Parametric Multi-Index Model,'' \emph{Advances in Water Resources}, vol. 57,
pp. 12--18, 2013.

\bibitem{friedman2001greedy}
J. H. Friedman, ``Greedy Function Approximation: A Gradient Boosting Machine,''
\emph{Annals of Statistics}, vol. 29, no. 5, pp. 1189--1232, 2001.

\end{thebibliography}

\end{document}
